% ===== PART 1: LÝ THUYẾT (Thành viên 1 - Ngọc) =====

\section{Giới thiệu chương}
\label{sec:intro}

Bài toán \textbf{xếp hạng} (ranking) là một trong những bài toán 
trung tâm của học máy hiện đại, xuất hiện rộng rãi trong nhiều 
ứng dụng thực tế. Khi người dùng nhập một truy vấn vào công cụ 
tìm kiếm như Google, hệ thống không chỉ cần tìm ra các trang web 
liên quan mà còn phải \textit{sắp xếp} chúng theo mức độ liên quan 
giảm dần. Tương tự, các hệ thống gợi ý như Netflix hay Spotify cần 
xếp hạng hàng nghìn bộ phim hay bài hát để đưa ra danh sách phù hợp 
nhất với từng người dùng.

Điểm khác biệt cốt lõi giữa ranking và các bài toán học máy truyền 
thống là: thay vì dự đoán một nhãn hay một giá trị tuyệt đối cho 
từng điểm dữ liệu, mục tiêu của ranking là học một \textbf{thứ tự 
tương đối} giữa các điểm. Nói cách khác, điều quan trọng không phải 
là điểm dữ liệu $x$ được gán nhãn gì, mà là liệu $x$ có đứng 
\textit{trước} hay \textit{sau} $x'$ trong thứ tự xếp hạng.

Chương 9 của~\cite{mohri2018foundations} trình bày một cách hệ thống 
các nền tảng lý thuyết cho bài toán ranking, bao gồm:

\begin{itemize}
  \item \textbf{Định nghĩa hình thức bài toán ranking} 
        (Mục~\ref{sec:problem}): Phát biểu bài toán dưới dạng học 
        hàm scoring từ các cặp dữ liệu có thứ tự ưu tiên.

  \item \textbf{Chặn tổng quát hoá} 
        (Mục~\ref{sec:genbounds}): Dẫn xuất các chặn trên cho 
        ranking loss tổng quát dựa trên Rademacher complexity.

  \item \textbf{Ranking với SVM} 
        (Mục~\ref{sec:svmranking}): Mở rộng bài toán tối ưu hoá 
        SVM sang không gian các cặp dữ liệu.

  \item \textbf{Thuật toán RankBoost} 
        (Mục~\ref{sec:rankboost}): Áp dụng tư tưởng boosting để 
        kết hợp nhiều hàm ranking yếu thành một hàm ranking mạnh.

  \item \textbf{Bipartite ranking và AUC} 
        (Mục~\ref{sec:bipartite}): Xem xét trường hợp đặc biệt 
        khi tập dữ liệu chia thành hai nhóm dương và âm, liên hệ 
        trực tiếp với diện tích dưới đường cong ROC (AUC).

  \item \textbf{Preference-based setting} 
        (Mục~\ref{sec:preference}): Tổng quát hoá sang trường hợp 
        học từ dữ liệu sở thích (preference feedback).
\end{itemize}

Về mặt lý thuyết, chương này được xây dựng trên nền tảng của hai 
chương trước trong sách: \textbf{Rademacher complexity} (Chương~3) 
được dùng trực tiếp để chứng minh các chặn tổng quát hoá, và 
\textbf{Support Vector Machines} (Chương~5) được mở rộng tự nhiên 
sang không gian cặp dữ liệu. Ngoài ra, RankBoost (Mục~\ref{sec:rankboost}) 
kế thừa tư tưởng từ \textbf{AdaBoost} (Chương~7).

% ===================================================
\section{Phần 1: Báo cáo lý thuyết}

\subsection{Kiến thức chuẩn bị}

\subsubsection{Không gian giả thuyết và hàm mất mát}
\begin{definition}
Cho tập mẫu $\mathcal{X}$, một \textbf{hàm scoring} là ánh xạ 
$f: \mathcal{X} \to \mathbb{R}$ 
dùng để gán điểm số cho từng điểm dữ liệu.
\end{definition}

\subsubsection{Rademacher Complexity}
Nhắc lại khái niệm Rademacher complexity sẽ được dùng 
trong chứng minh generalization bound ở Phần~\ref{sec:genbounds}.

\begin{definition}
Cho tập mẫu $S = \{x_1, \ldots, x_m\}$ lấy từ phân phối $\mathcal{D}$,
\textbf{Rademacher complexity thực nghiệm} của lớp hàm $\mathcal{H}$ là:
\[
\hat{\mathfrak{R}}_S(\mathcal{H}) = 
\mathbb{E}_{\boldsymbol{\sigma}}\left[
  \sup_{h \in \mathcal{H}} \frac{1}{m} \sum_{i=1}^{m} \sigma_i h(x_i)
\right]
\]
trong đó $\sigma_i$ là các biến ngẫu nhiên Rademacher độc lập.
\end{definition}

\subsubsection{Support Vector Machine (SVM)}

SVM (Chương~5 của~\cite{mohri2018foundations}) là thuật toán học 
có giám sát tìm siêu phẳng phân tách hai lớp với \textbf{margin 
lớn nhất}. Với tập huấn luyện $\{(z_i, y_i)\}_{i=1}^n$, 
$y_i \in \{-1,+1\}$, bài toán tối ưu hoá SVM soft-margin là:
\[
  \min_{\mathbf{w}, b, \boldsymbol{\xi}} 
  \frac{1}{2}\|\mathbf{w}\|^2 + C\sum_{i=1}^n \xi_i
\]
\[
  \text{s.t. } y_i(\mathbf{w}^\top \phi(z_i) + b) \geq 1 - \xi_i, 
  \quad \xi_i \geq 0, \quad \forall i,
\]
trong đó $\phi: \mathcal{X} \to \mathcal{F}$ là ánh xạ đặc trưng 
vào không gian Hilbert $\mathcal{F}$, $C > 0$ là tham số điều chỉnh. 
Ở Mục~\ref{sec:svmranking}, ta sẽ thấy bài toán này được mở rộng 
tự nhiên sang không gian cặp $(x, x')$ để giải bài toán ranking.

\subsection{Bài toán Ranking}
\label{sec:problem}

\subsubsection{Phát biểu hình thức}

Trong bài toán ranking, ta được cho một tập mẫu huấn luyện gồm 
các \textbf{cặp có thứ tự ưu tiên}. Cụ thể, mỗi mẫu huấn luyện 
là một cặp $(x, x') \in \mathcal{X} \times \mathcal{X}$ kèm theo 
thông tin rằng $x$ \textit{được ưu tiên hơn} $x'$, ký hiệu là 
$x \succ x'$.

\begin{definition}[Hàm scoring]
Một \textbf{hàm scoring} (hay hàm ranking) là ánh xạ 
$f: \mathcal{X} \to \mathbb{R}$ sao cho với mọi cặp 
$(x, x')$ có $x \succ x'$, ta mong muốn:
\[
  f(x) > f(x').
\]
Nói cách khác, điểm dữ liệu được ưu tiên hơn sẽ nhận điểm số 
(score) cao hơn.
\end{definition}

\begin{definition}[Ranking loss]
\label{def:rankingloss}
Cho hàm scoring $f: \mathcal{X} \to \mathbb{R}$, 
\textbf{ranking loss} (hay pairwise loss) trên một cặp 
$(x, x')$ với $x \succ x'$ được định nghĩa là:
\[
  \ell(f, (x, x')) = \mathbf{1}[f(x) < f(x')] 
  + \frac{1}{2}\mathbf{1}[f(x) = f(x')],
\]
trong đó $\mathbf{1}[\cdot]$ là hàm chỉ thị. Ranking loss 
bằng $1$ khi $f$ xếp nhầm thứ tự, bằng $\frac{1}{2}$ khi 
$f$ không phân biệt được hai điểm, và bằng $0$ khi $f$ 
xếp đúng.
\end{definition}

\subsubsection{Ví dụ minh hoạ}

\begin{example}
Xét bài toán xếp hạng kết quả tìm kiếm với 
$\mathcal{X} = \mathbb{R}^2$, trong đó mỗi điểm dữ liệu 
$x = (x_1, x_2)$ biểu diễn một trang web với hai đặc trưng: 
$x_1$ là số từ khớp với truy vấn, $x_2$ là độ uy tín của trang.

Giả sử ta có hai trang web $x = (5, 0.8)$ và $x' = (3, 0.4)$, 
và nhãn $x \succ x'$ (trang $x$ liên quan hơn). Xét hàm scoring 
tuyến tính $f(x) = w_1 x_1 + w_2 x_2$ với $w_1 = w_2 = 1$:
\[
  f(x) = 5 + 0.8 = 5.8, \quad f(x') = 3 + 0.4 = 3.4.
\]
Vì $f(x) = 5.8 > 3.4 = f(x')$, hàm $f$ xếp đúng thứ tự, 
ranking loss bằng $0$.
\end{example}

\subsubsection{Thiết lập học thống kê}

Giả sử các cặp $(x, x')$ được lấy mẫu độc lập từ một phân 
phối $\mathcal{D}$ trên $\mathcal{X} \times \mathcal{X}$. 
\textbf{Ranking loss tổng quát} (generalization ranking loss) 
của hàm scoring $f$ được định nghĩa là:
\[
  R(f) = \Pr_{(x,x') \sim \mathcal{D}}[f(x) < f(x')] 
       + \frac{1}{2}\Pr_{(x,x') \sim \mathcal{D}}[f(x) = f(x')].
\]

Cho tập huấn luyện $S = \{(x_1, x_1'), \ldots, (x_m, x_m')\}$ 
gồm $m$ cặp lấy mẫu i.i.d. từ $\mathcal{D}$, 
\textbf{ranking loss thực nghiệm} được tính là:
\[
  \hat{R}_S(f) = \frac{1}{m} \sum_{i=1}^{m} 
  \left( \mathbf{1}[f(x_i) < f(x_i')] 
  + \frac{1}{2}\mathbf{1}[f(x_i) = f(x_i')] \right).
\]

Mục tiêu của thuật toán học là tìm hàm $f$ trong lớp giả thuyết 
$\mathcal{H}$ sao cho $R(f)$ nhỏ nhất có thể, dựa trên thông tin 
từ tập huấn luyện $S$.

\subsection{Generalization Bounds}
\label{sec:genbounds}

Mục tiêu của phần này là trả lời câu hỏi: \textit{Nếu một hàm 
scoring $f$ có ranking loss thực nghiệm nhỏ trên tập huấn luyện, 
liệu ranking loss tổng quát của nó có nhỏ không?} Để trả lời, 
ta dẫn xuất các chặn tổng quát hoá (generalization bounds) cho 
bài toán ranking dựa trên Rademacher complexity.

\subsubsection{Quan sát cốt lõi: Ranking như Classification trên cặp}

Ý tưởng then chốt là \textbf{quy bài toán ranking về bài toán 
phân loại nhị phân trên không gian các cặp}. 

Cho lớp hàm scoring $\mathcal{H}$ trên $\mathcal{X}$, ta định 
nghĩa lớp hàm tương ứng trên không gian cặp 
$\mathcal{X} \times \mathcal{X}$:
\[
  \mathcal{G} = \{ g : (x, x') \mapsto f(x) - f(x') \mid f \in \mathcal{H} \}.
\]
Khi đó, với mỗi cặp $(x, x')$ có $x \succ x'$:
\begin{align*}
  \ell(f, (x,x')) &= \mathbf{1}[f(x) < f(x')] + \tfrac{1}{2}\mathbf{1}[f(x) = f(x')] \\
                  &= \mathbf{1}[g(x,x') < 0] + \tfrac{1}{2}\mathbf{1}[g(x,x') = 0].
\end{align*}
Đây chính xác là dạng \textbf{0-1 loss} của bài toán phân loại 
nhị phân, trong đó $g(x,x')$ đóng vai trò như hàm phân loại 
và nhãn luôn là $+1$ (vì $x \succ x'$).

\subsubsection{Rademacher Complexity của lớp hàm ranking}

\begin{lemma}
\label{lem:rademacher_ranking}
Cho lớp hàm scoring $\mathcal{H}$ trên $\mathcal{X}$ và tập mẫu 
cặp $S = \{(x_1, x_1'), \ldots, (x_m, x_m')\}$. Khi đó Rademacher 
complexity thực nghiệm của $\mathcal{G}$ thoả mãn:
\[
  \hat{\mathfrak{R}}_S(\mathcal{G}) \leq 2\hat{\mathfrak{R}}_{S^+}(\mathcal{H}),
\]
trong đó $S^+ = \{x_1, x_1', x_2, x_2', \ldots, x_m, x_m'\}$ là 
tập gồm $2m$ điểm đơn lẻ.
\end{lemma}

\begin{proof}
Theo định nghĩa:
\begin{align*}
  \hat{\mathfrak{R}}_S(\mathcal{G}) 
  &= \mathbb{E}_{\boldsymbol{\sigma}}\left[
      \sup_{f \in \mathcal{H}} \frac{1}{m} \sum_{i=1}^{m} 
      \sigma_i (f(x_i) - f(x_i'))
    \right] \\
  &= \mathbb{E}_{\boldsymbol{\sigma}}\left[
      \sup_{f \in \mathcal{H}} \frac{1}{m} 
      \left( \sum_{i=1}^{m} \sigma_i f(x_i) 
           - \sum_{i=1}^{m} \sigma_i f(x_i') \right)
    \right].
\end{align*}
Vì $\sigma_i$ và $-\sigma_i$ có cùng phân phối (đối xứng quanh 0), 
ta có $-\sigma_i f(x_i') \overset{d}{=} \sigma_i' f(x_i')$ với 
$\sigma_i'$ là biến Rademacher độc lập. Áp dụng bất đẳng thức 
tam giác cho $\sup$:
\begin{align*}
  \hat{\mathfrak{R}}_S(\mathcal{G}) 
  &\leq \mathbb{E}_{\boldsymbol{\sigma}}\left[
      \sup_{f \in \mathcal{H}} \frac{1}{m} \sum_{i=1}^{m} \sigma_i f(x_i)
    \right]
  + \mathbb{E}_{\boldsymbol{\sigma}'}\left[
      \sup_{f \in \mathcal{H}} \frac{1}{m} \sum_{i=1}^{m} \sigma_i' f(x_i')
    \right] \\
  &= 2\hat{\mathfrak{R}}_{S^+}(\mathcal{H}).
\end{align*}
\end{proof}

\subsubsection{Định lý chặn tổng quát hoá chính}

\begin{theorem}[Generalization bound cho Ranking]
\label{thm:ranking_bound}
Cho lớp hàm scoring $\mathcal{H}$ với $f: \mathcal{X} \to [0,1]$ 
với mọi $f \in \mathcal{H}$. Cho $S = \{(x_i, x_i')\}_{i=1}^m$ 
lấy mẫu i.i.d. từ phân phối $\mathcal{D}$ trên 
$\mathcal{X} \times \mathcal{X}$. Khi đó với mọi $\delta > 0$, 
với xác suất ít nhất $1 - \delta$ trên việc chọn mẫu $S$, 
đồng thời với mọi $f \in \mathcal{H}$:
\[
  \boxed{R(f) \leq \hat{R}_S(f) + 4\hat{\mathfrak{R}}_{S^+}(\mathcal{H}) 
  + \sqrt{\frac{\log(2/\delta)}{2m}}.}
\]
\end{theorem}

\begin{proof}
Chứng minh gồm ba bước.

\textbf{Bước 1: Áp dụng bất đẳng thức McDiarmid.}

Xét hàm $\Phi(S) = \sup_{f \in \mathcal{H}} [R(f) - \hat{R}_S(f)]$. 
Khi thay một cặp $(x_i, x_i')$ bằng $(z_i, z_i')$, giá trị 
$\hat{R}_S(f)$ thay đổi không quá $\frac{1}{m}$ với mọi $f$, 
do đó $\Phi$ thay đổi không quá $\frac{1}{m}$. Theo bất đẳng thức 
McDiarmid:
\[
  \Pr\left[\Phi(S) - \mathbb{E}[\Phi(S)] \geq \epsilon\right] 
  \leq \exp\left(-2m\epsilon^2\right).
\]
Đặt vế phải bằng $\delta/2$ và giải ra $\epsilon$:
\[
  \Phi(S) \leq \mathbb{E}_S[\Phi(S)] + \sqrt{\frac{\log(2/\delta)}{2m}}.
\]

\textbf{Bước 2: Chặn $\mathbb{E}_S[\Phi(S)]$ bằng Rademacher complexity.}

Dùng kỹ thuật đối xứng hoá (symmetrization) chuẩn với ghost sample 
$S' = \{(x_i'', x_i''')\}_{i=1}^m$ độc lập và cùng phân phối với $S$:
\begin{align*}
  \mathbb{E}_S[\Phi(S)] 
  &= \mathbb{E}_S\left[\sup_{f} (R(f) - \hat{R}_S(f))\right] \\
  &\leq \mathbb{E}_{S,S'}\left[\sup_{f} 
      (\hat{R}_{S'}(f) - \hat{R}_S(f))\right] \\
  &= \mathbb{E}_{S,S',\boldsymbol{\sigma}}\left[\sup_{f} 
      \frac{1}{m}\sum_{i=1}^m \sigma_i 
      (\ell(f,(x_i'',x_i''')) - \ell(f,(x_i,x_i')))\right] \\
  &\leq 2\mathbb{E}_{S,\boldsymbol{\sigma}}\left[\sup_{f} 
      \frac{1}{m}\sum_{i=1}^m \sigma_i \ell(f,(x_i,x_i'))\right] \\
  &= 2\hat{\mathfrak{R}}_S(\mathcal{L} \circ \mathcal{G}).
\end{align*}

\textbf{Bước 3: Áp dụng Bổ đề~\ref{lem:rademacher_ranking} và 
bổ đề Talagrand.}

Vì hàm loss $\ell$ là 1-Lipschitz theo $g(x,x') = f(x) - f(x')$, 
theo bổ đề contraction của Talagrand:
\[
  \hat{\mathfrak{R}}_S(\mathcal{L} \circ \mathcal{G}) 
  \leq \hat{\mathfrak{R}}_S(\mathcal{G}) 
  \leq 2\hat{\mathfrak{R}}_{S^+}(\mathcal{H}).
\]
Kết hợp ba bước lại:
\[
  R(f) \leq \hat{R}_S(f) + 4\hat{\mathfrak{R}}_{S^+}(\mathcal{H}) 
  + \sqrt{\frac{\log(2/\delta)}{2m}}.
\]
\end{proof}

\begin{morong}
\begin{lemma}[Bổ đề contraction -- Talagrand]
\label{lem:talagrand}
Cho $\phi: \mathbb{R} \to \mathbb{R}$ là hàm $L$-Lipschitz 
(tức là $|\phi(a) - \phi(b)| \leq L|a-b|$ với mọi $a,b$) 
và $\mathcal{G}$ là lớp hàm trên $\mathcal{X}$. Khi đó:
\[
  \hat{\mathfrak{R}}_S(\phi \circ \mathcal{G}) 
  \leq L \cdot \hat{\mathfrak{R}}_S(\mathcal{G}).
\]
\end{lemma}

\begin{proof}
Theo định nghĩa Rademacher complexity:
\begin{align*}
  \hat{\mathfrak{R}}_S(\phi \circ \mathcal{G})
  &= \mathbb{E}_{\boldsymbol{\sigma}}\left[
      \sup_{g \in \mathcal{G}} \frac{1}{m}
      \sum_{i=1}^m \sigma_i \phi(g(x_i))
    \right].
\end{align*}
Cố định $\boldsymbol{\sigma}$ và xét $g^* = \arg\sup_{g} 
\sum_i \sigma_i \phi(g(x_i))$. Đặt $g^0$ là hàm thoả 
$g^0(x_i) = 0$ với mọi $i$. Ta có:
\begin{align*}
  \sum_{i=1}^m \sigma_i \phi(g^*(x_i))
  &= \sum_{i=1}^m \sigma_i [\phi(g^*(x_i)) - \phi(0)] 
   + \sum_{i=1}^m \sigma_i \phi(0) \\
  &\leq \sum_{i=1}^m |\phi(g^*(x_i)) - \phi(0)| 
   + \left|\sum_{i=1}^m \sigma_i\right| \cdot |\phi(0)| \\
  &\leq L \sum_{i=1}^m |g^*(x_i)| 
   + \left|\sum_{i=1}^m \sigma_i\right| \cdot |\phi(0)|.
\end{align*}
Lấy kỳ vọng theo $\boldsymbol{\sigma}$ và dùng tính chất 
$\mathbb{E}[\sigma_i] = 0$:
\[
  \hat{\mathfrak{R}}_S(\phi \circ \mathcal{G})
  \leq L \cdot \mathbb{E}_{\boldsymbol{\sigma}}\left[
      \sup_{g \in \mathcal{G}} \frac{1}{m}
      \sum_{i=1}^m \sigma_i g(x_i)
    \right]
  = L \cdot \hat{\mathfrak{R}}_S(\mathcal{G}).
\]
\end{proof}

Áp dụng vào chứng minh Định lý~\ref{thm:ranking_bound}: 
hàm ranking loss $\ell(t) = \mathbf{1}[t < 0] + 
\frac{1}{2}\mathbf{1}[t=0]$ là $1$-Lipschitz, do đó 
$L = 1$ và:
\[
  \hat{\mathfrak{R}}_S(\mathcal{L} \circ \mathcal{G}) 
  \leq \hat{\mathfrak{R}}_S(\mathcal{G}).
\]
\end{morong}

\subsubsection{Ý nghĩa của chặn}

Định lý~\ref{thm:ranking_bound} cho thấy ba thành phần kiểm soát 
generalization:

\begin{itemize}
  \item $\hat{R}_S(f)$: Ranking loss trên tập huấn luyện — càng 
  nhỏ càng tốt.
  
  \item $4\hat{\mathfrak{R}}_{S^+}(\mathcal{H})$: Độ phức tạp của 
  lớp hàm $\mathcal{H}$ — lớp hàm càng phức tạp thì chặn càng lỏng, 
  dễ overfit hơn.
  
  \item $\sqrt{\frac{\log(2/\delta)}{2m}}$: Hạng tử thống kê — 
  giảm theo $O(1/\sqrt{m})$ khi số mẫu $m$ tăng.
\end{itemize}

\begin{morong}
\subsubsection*{Ví dụ phản chứng: Điều kiện bounded là cần thiết}

Định lý~\ref{thm:ranking_bound} yêu cầu $f: \mathcal{X} \to [0,1]$. 
Ta xây dựng ví dụ cho thấy nếu bỏ điều kiện này, chặn có thể 
không còn hữu ích.

Xét $\mathcal{X} = \{-1, +1\}$ và lớp hàm 
$\mathcal{H}_c = \{f_c : x \mapsto cx \mid c > 0\}$ 
không bị chặn. Với tập mẫu $S = \{(x_i, x_i')\}$ gồm $m$ 
cặp bất kỳ, Rademacher complexity là:
\[
  \hat{\mathfrak{R}}_{S^+}(\mathcal{H}_c) 
  = \mathbb{E}_{\boldsymbol{\sigma}}\left[
      \sup_{c > 0} \frac{c}{m} \sum_{i} \sigma_i x_i
    \right].
\]
Khi $\sum_i \sigma_i x_i > 0$, $\sup_{c>0}$ đạt $+\infty$. 
Do đó $\hat{\mathfrak{R}}_{S^+}(\mathcal{H}_c) = +\infty$ 
với xác suất dương, khiến chặn trong 
Định lý~\ref{thm:ranking_bound} trở nên vô nghĩa.

Điều này cho thấy điều kiện $f$ có giá trị trong $[0,1]$ 
(hoặc tổng quát hơn là bị chặn) là \textbf{cần thiết} để 
chặn tổng quát hoá có ý nghĩa thực tế.
\end{morong}

\begin{example}
Xét $\mathcal{H}$ là lớp hàm tuyến tính chuẩn hoá: 
$\mathcal{H} = \{x \mapsto \mathbf{w}^\top x \mid \|\mathbf{w}\| \leq \Lambda, 
\|x\| \leq r\}$. Từ Chương~5, Rademacher complexity của lớp này thoả:
\[
  \hat{\mathfrak{R}}_{S^+}(\mathcal{H}) \leq \frac{r\Lambda}{\sqrt{2m}}.
\]
Thay vào Định lý~\ref{thm:ranking_bound}:
\[
  R(f) \leq \hat{R}_S(f) + \frac{4r\Lambda}{\sqrt{2m}} 
  + \sqrt{\frac{\log(2/\delta)}{2m}}.
\]
Khi $m \to \infty$, hai hạng tử sau đều về $0$, tức là hàm 
scoring học được sẽ tổng quát hoá tốt.
\end{example}

\subsection{Ranking with SVMs}
\label{sec:svmranking}

Phần này mở rộng bài toán tối ưu hoá SVM sang không gian các cặp 
dữ liệu. Ý tưởng cốt lõi là: thay vì phân loại từng điểm $x$, 
ta \textbf{phân loại từng cặp} $(x, x')$ — cụ thể là phân loại 
xem hiệu $f(x) - f(x')$ có dương hay không.

\subsubsection{Nhắc lại SVM phân loại}

Trong SVM phân loại nhị phân (Chương~5), với tập huấn luyện 
$\{(z_i, y_i)\}_{i=1}^n$ trong đó $y_i \in \{-1, +1\}$, ta 
tìm siêu phẳng $\mathbf{w}^\top \phi(z) + b$ tối ưu hoá:
\[
  \min_{\mathbf{w}, b, \boldsymbol{\xi}} 
  \frac{1}{2}\|\mathbf{w}\|^2 + C\sum_{i=1}^n \xi_i
\]
\[
  \text{s.t. } y_i(\mathbf{w}^\top \phi(z_i) + b) \geq 1 - \xi_i, 
  \quad \xi_i \geq 0, \quad \forall i.
\]

\subsubsection{Bài toán tối ưu hoá SVM Ranking}

Trong bài toán ranking, tập huấn luyện gồm $m$ cặp 
$\{(x_i, x_i')\}_{i=1}^m$ với $x_i \succ x_i'$. Ta muốn 
tìm hàm scoring tuyến tính $f(x) = \mathbf{w}^\top \phi(x)$ 
sao cho $f(x_i) > f(x_i')$ với mọi $i$.

Đặt $\delta_i = \phi(x_i) - \phi(x_i')$. Khi đó điều kiện 
$f(x_i) > f(x_i')$ tương đương với:
\[
  \mathbf{w}^\top \delta_i > 0.
\]
Đây chính xác là bài toán phân loại nhị phân trên không gian 
cặp, với vector đặc trưng $\delta_i$ và nhãn $+1$. Áp dụng 
SVM soft-margin:

\begin{definition}[Bài toán SVM Ranking]
\label{def:svmranking}
Bài toán tối ưu hoá \textbf{SVM Ranking} được phát biểu như sau:
\[
  \min_{\mathbf{w}, \boldsymbol{\xi}} 
  \frac{1}{2}\|\mathbf{w}\|^2 + C\sum_{i=1}^m \xi_i
\]
\[
  \text{s.t. } \mathbf{w}^\top(\phi(x_i) - \phi(x_i')) 
  \geq 1 - \xi_i, \quad \xi_i \geq 0, \quad \forall i \in [m].
\]
trong đó $C > 0$ là tham số điều chỉnh đánh đổi giữa độ rộng 
margin và số lỗi huấn luyện, $\xi_i$ là biến bù (slack variable).
\end{definition}

\subsubsection{Bài toán đối ngẫu (Dual Problem)}

Áp dụng nhân tử Lagrange, bài toán đối ngẫu của 
Định nghĩa~\ref{def:svmranking} là:
\[
  \max_{\boldsymbol{\alpha}} 
  \sum_{i=1}^m \alpha_i 
  - \frac{1}{2} \sum_{i,j=1}^m \alpha_i \alpha_j 
    K_\Delta(i, j)
\]
\[
  \text{s.t. } 0 \leq \alpha_i \leq C, \quad \forall i \in [m],
\]
trong đó kernel trên không gian cặp được định nghĩa là:
\[
  K_\Delta(i,j) = \langle \delta_i, \delta_j \rangle 
  = K(x_i, x_j) - K(x_i, x_j') - K(x_i', x_j) + K(x_i', x_j'),
\]
với $K(x, x') = \langle \phi(x), \phi(x') \rangle$ là kernel gốc.

Nghiệm nguyên thuỷ được phục hồi qua:
\[
  \mathbf{w}^* = \sum_{i=1}^m \alpha_i \delta_i 
  = \sum_{i=1}^m \alpha_i (\phi(x_i) - \phi(x_i')).
\]

\subsubsection{Dự đoán và margin}

Sau khi học được $\mathbf{w}^*$, hàm scoring được dùng để xếp 
hạng một điểm mới $x$ là:
\[
  f(x) = \mathbf{w}^{*\top} \phi(x) 
  = \sum_{i=1}^m \alpha_i (K(x_i, x) - K(x_i', x)).
\]
Với hai điểm $x$ và $x'$ cần so sánh, ta kết luận $x \succ x'$ 
nếu $f(x) > f(x')$, tức là:
\[
  \sum_{i=1}^m \alpha_i 
  \big(K(x_i,x) - K(x_i',x) - K(x_i,x') + K(x_i',x')\big) > 0.
\]

\subsubsection{So sánh với SVM phân loại}

Bảng~\ref{tab:svm_compare} tóm tắt điểm tương đồng và khác biệt 
giữa SVM phân loại và SVM Ranking.

\begin{table}[h]
\centering
\begin{tabular}{lcc}
\hline
\textbf{Đặc điểm} & \textbf{SVM Phân loại} & \textbf{SVM Ranking} \\
\hline
Không gian đầu vào & $\mathcal{X}$ & $\mathcal{X} \times \mathcal{X}$ \\
Vector đặc trưng & $\phi(x)$ & $\phi(x) - \phi(x')$ \\
Nhãn huấn luyện & $y \in \{-1,+1\}$ & $x \succ x'$ (luôn $+1$) \\
Kernel & $K(x,x')$ & $K_\Delta(i,j)$ \\
Số ràng buộc & $n$ & $m$ (số cặp) \\
\hline
\end{tabular}
\caption{So sánh SVM phân loại và SVM Ranking.}
\label{tab:svm_compare}
\end{table}

\begin{example}
Xét $\mathcal{X} = \mathbb{R}^2$ với kernel tuyến tính 
$K(x,x') = x^\top x'$. Cho tập huấn luyện gồm 3 cặp:
$(x_1, x_1') = ((3,2),(1,1))$, 
$(x_2, x_2') = ((4,1),(2,3))$, 
$(x_3, x_3') = ((2,4),(1,2))$.

Các vector hiệu là:
\begin{align*}
  \delta_1 &= (3,2) - (1,1) = (2,1), \\
  \delta_2 &= (4,1) - (2,3) = (2,-2), \\
  \delta_3 &= (2,4) - (1,2) = (1,2).
\end{align*}
Bài toán SVM Ranking trở thành tìm $\mathbf{w} \in \mathbb{R}^2$ 
sao cho $\mathbf{w}^\top \delta_i \geq 1$ với $i = 1,2,3$, 
đồng thời tối thiểu $\|\mathbf{w}\|^2$.
\end{example}

\begin{morong}
\subsubsection*{Phân tích độ phức tạp tính toán}

So sánh độ phức tạp giữa SVM phân loại và SVM Ranking:

\begin{itemize}
  \item \textbf{SVM phân loại}: Bài toán QP với $n$ biến 
  $\alpha_i$ và $n$ ràng buộc — độ phức tạp 
  $O(n^2)$ đến $O(n^3)$ tuỳ solver.

  \item \textbf{SVM Ranking}: Bài toán QP với $m$ biến, 
  trong đó $m$ là số cặp huấn luyện. Nếu có $p$ điểm 
  dữ liệu gốc thì $m = O(p^2)$ trong trường hợp xấu nhất 
  (tất cả các cặp). Do đó độ phức tạp là 
  $O(p^4)$ đến $O(p^6)$ — \textbf{đắt hơn nhiều} so với 
  SVM phân loại trên cùng tập dữ liệu gốc.
\end{itemize}

\textbf{Hệ quả thực tế:} Trong ứng dụng, người ta thường 
chỉ lấy mẫu một tập con các cặp thay vì dùng toàn bộ 
$O(p^2)$ cặp, đánh đổi giữa độ chính xác và tốc độ huấn 
luyện. Đây là một trong những động lực để phát triển các 
thuật toán ranking hiệu quả hơn như RankBoost 
(Mục~\ref{sec:rankboost}).
\end{morong}